% =============================================================================
% CORRECTIONS FILE
% Replace the indicated blocks in the main .tex with the versions below.
% Everything is 100% identical to the original except the marked fixes.
% =============================================================================


% =============================================================================
% REPLACE: \subsection{Selection Schemes} block  (was lines 263-277)
% FIXES: elitism top-2 -> top-1 | survival scheme description
% =============================================================================

We evaluate three selection configurations:

\begin{enumerate}
    \item \textbf{Tournament selection + crossover + mutation}: Three individuals are drawn
    at random; the highest-fitness individual is admitted to the mating pool.

    \item \textbf{Elitism (top-1) + crossover + mutation}: The highest-fitness
    prompt is carried forward unchanged into every subsequent generation.

    \item \textbf{Roulette selection + crossover + mutation}: Each individual is selected
    with probability proportional to its fitness score.
\end{enumerate}

\textbf{Survival scheme:} At each generation all non-elite individuals are replaced by
newly bred offspring; the population is fully regenerated except for the elite
carry-forward (one individual in the elitism scheme, zero in the tournament and roulette
schemes).


% =============================================================================
% REPLACE: \textbf{Crossover} sentence in \subsection{Evolutionary Operators}
%          (was lines 293-297)
% FIX: quoted meta-prompt matches actual code in llm.py
% =============================================================================

\textbf{Crossover (rate 0.8):} Given two parent prompts $P_A$ and $P_B$, a meta-prompt is
submitted to the LLM: \textit{``Combine these two coding-instruction prompts into one
coherent prompt. Preserve the strongest strategies from each parent. Do not simply
concatenate --- produce a unified instruction paragraph.''} This produces an offspring
that inherits strategic elements from both parents while remaining grammatically coherent.


% =============================================================================
% REPLACE: Table 2 (Experimental Parameters) tabular body  (was lines 354-367)
% FIXES: Survival pruning row | Target LLM | add Optimizer LLM | temperatures
% =============================================================================

\begin{tabular}{ll}
\toprule
\textbf{Parameter} & \textbf{Value} \\
\midrule
Population size & 5 \\
Training problems & 65 \\
Generations & 12 \\
Crossover rate & 0.8 \\
Mutation: strategy injection & 0.25 \\
Mutation: sentence deletion & 0.15 \\
Mutation: rephrasing & 0.10 \\
Survival & Full generational replacement \\
Total API calls (approx.) & 328--334 per generation \\
Target LLM & Gemini 2.5 Flash (Google, Vertex AI) \\
Optimizer LLM & Gemini 2.5 Pro (Google, Vertex AI) \\
LLM temperature & 0.0 (evaluation), 0.7 (operators) \\
\bottomrule
\end{tabular}


% =============================================================================
% REPLACE: entire \section{Results} block  (was lines 370-432)
% FIX: \label commands now match the \ref{} keys used in the Discussion section
% =============================================================================

%=============================================================================
\section{Results}
\label{sec:results}
%=============================================================================

We present results for each of the three selection schemes, followed by an overall
comparison. Each figure shows best, mean, and worst fitness (pass@1) per generation.

\subsection{Roulette Selection}
\begin{figure}[H]
    \centering
    \includegraphics[width=\linewidth]{roulette_convergence.png}
    \caption{Convergence}
    \label{fig:roulette}
\end{figure}
\begin{figure}[H]
    \centering
    \includegraphics[width=\linewidth]{roulette_diversity.png}
    \caption{Diversity}
    \label{fig:roulette_div}
\end{figure}

\subsection{Tournament Selection with Elitism}
\begin{figure}[H]
    \centering
    \includegraphics[width=\linewidth]{tournament_convergence.png}
    \caption{Convergence}
    \label{fig:tourn_elite}
\end{figure}
\begin{figure}[H]
    \centering
    \includegraphics[width=\linewidth]{tournament_diversity.png}
    \caption{Diversity}
    \label{fig:tourn_elite_div}
\end{figure}

\subsection{Tournament Selection (No Elitism)}
\begin{figure}[H]
    \centering
    \includegraphics[width=\linewidth]{tournament_convergence.png}
    \caption{Convergence}
    \label{fig:tourn}
\end{figure}
\begin{figure}[H]
    \centering
    \includegraphics[width=\linewidth]{tournament_diversity.png}
    \caption{Diversity}
    \label{fig:tourn_div}
\end{figure}

\subsection{Overall Comparison}
\begin{figure}[H]
    \centering
    \includegraphics[width=\linewidth]{test_3_score_by_rating.png}
    \caption{Score with respect to different difficulty rating}
    \label{fig:overall}
\end{figure}
\begin{figure}[H]
    \centering
    \includegraphics[width=\linewidth]{test_4_per_problem_scatter.png}
    \caption{Handwritten Prompts Vs Evolved Prompts}
    \label{fig:seed_vs_evolved}
\end{figure}
